\section*{How the Proof Works}

For the global three-dimensional Navier--Stokes regularity problem, the fluid
obeys
\[
\partial_tu+(u\!\cdot\nabla)u+\nabla p=\nu\Delta u,
\qquad
\nabla\!\cdot u=0,
\qquad
u(\cdot,0)=u_0.
\]
Here \(u\) is the velocity, \(p\) is the pressure fixed by incompressibility,
and \(\nu>0\) is the viscosity. The nonlinear term moves and deforms the
velocity field. Viscosity diffuses its spatial variation.

For a smooth solution that decays at infinity, taking the \(L^2\) inner product
with \(u\) gives
\[
\frac12\|u(t)\|_2^2
+\nu\int_0^t\|\nabla u(\tau)\|_2^2\,d\tau
=\frac12\|u_0\|_2^2.
\]
This bounds the kinetic energy at each time and the accumulated gradient over
the interval. It does not bound \(\|\nabla u(t)\|_2\) at each individual time.

As a vortex tube stretches along its axis, the parcels of rotating fluid
inside it lengthen in that direction. Incompressibility preserves each
parcel's volume, so the axial extension is balanced by contraction across the
tube. The rotation is now carried by a thinner core. With
\(\omega=\nabla\times u\), taking the curl of the equation gives
\[
D_t\omega
:=\partial_t\omega+(u\!\cdot\nabla)\omega
=(\omega\!\cdot\nabla)u+\nu\Delta\omega.
\]
The term \((\omega\!\cdot\nabla)u\) is vortex stretching. When the surrounding
strain pulls in the direction of \(\omega\), this term increases
\(|\omega|\): the vorticity becomes stronger as the tube grows longer and
thinner. Since \(\omega\) is built from first spatial derivatives of \(u\),
the velocity changes more sharply across the core. Viscosity acts against that
sharpening through \(\nu\Delta\omega\).

To follow this balance into a smaller core, we need a precise way to shrink the
solution without changing the equation it solves. Navier--Stokes scaling builds
a narrower, faster field from \((u,p)\):
\[
u_\lambda(x,t)=\lambda u(\lambda x,\lambda^2t),
\qquad
p_\lambda(x,t)=\lambda^2p(\lambda x,\lambda^2t),
\qquad
\lambda>1.
\]
If the original structure has width \(r\), velocity scale \(U\), and duration
\(\tau\), its rescaled copy has width \(r/\lambda\), velocity scale
\(\lambda U\), and duration \(\tau/\lambda^2\). Substitution gives
\[
\begin{aligned}
&\partial_tu_\lambda
+(u_\lambda\!\cdot\nabla)u_\lambda
+\nabla p_\lambda-\nu\Delta u_\lambda\\
&\qquad
=\lambda^3
\bigl[\partial_tu+(u\!\cdot\nabla)u+\nabla p-\nu\Delta u\bigr]
(\lambda x,\lambda^2t)
=0.
\end{aligned}
\]
The common factor \(\lambda^3\) cancels, so the rescaled fields obey the
original equation. Narrowing the core strengthens nonlinear sharpening and
viscous diffusion by an identical factor; scale alone gives viscosity no new
advantage.

The narrower field has velocity scale \(\lambda U\), but it occupies
\(\lambda^{-3}\) times the volume. Its kinetic energy scales as
\[
\|u_\lambda(\cdot,t)\|_2^2
=\lambda^{-1}\|u(\cdot,\lambda^2t)\|_2^2.
\]
The energy norm becomes smaller as the core becomes thinner, even while the
gradients become larger.

\begin{minipage}{\textwidth}
The proof needs a norm that does not shrink as the core narrows. For the
homogeneous Sobolev norms,
\[
\widehat{u_\lambda}(\xi,t)
=\lambda^{-2}\widehat u(\xi/\lambda,\lambda^2t),
\]
and a change of variables gives
\[
\begin{aligned}
\|u_\lambda(\cdot,t)\|_{\dot H^s}^2
&=\int_{\mathbb R^3}|\xi|^{2s}\lambda^{-4}
  |\widehat u(\xi/\lambda,\lambda^2t)|^2\,d\xi\\
&=\lambda^{2s-1}
  \|u(\cdot,\lambda^2t)\|_{\dot H^s}^2.
\end{aligned}
\]
At \(s=\tfrac12\), the factor is \(\lambda^0=1\). Contracting the field changes
its spatial scale without changing its \(\dot H^{1/2}\) size. Below that
height, concentration makes the norm smaller; above it, concentration makes
the norm larger. The energy norm becomes smaller under concentration; the
critical norm does not. That is why the proof begins at \(s=\tfrac12\).
\end{minipage}

The factor \(\lambda^2\) in time tells us how long each smaller turn would
take. Suppose one evolving vortex repeats the full Navier--Stokes rescaling
from turn to turn, with its width contracted by a fixed factor
\(0<\rho<1\). Then
\[
r_j=\rho^jr_0,
\qquad
U_j=\rho^{-j}U_0,
\qquad
\tau_j:=\frac{r_j}{U_j}=\rho^{2j}\tau_0.
\]
The duration of each turn is \(\rho^2\) times the preceding turn, and
\[
\sum_{j=0}^{\infty}\tau_j
=\tau_0\sum_{j=0}^{\infty}\rho^{2j}
=\frac{\tau_0}{1-\rho^2}<\infty.
\]
Under this repeated-contraction hypothesis, infinitely many turns fit before
a finite time while \(U_j\) diverges. This is the Zeno cascade: a thinner and
faster vortex at every turn, with nonlinear sharpening and viscous diffusion
still scaled together. Finite energy and finite duration do not, by
themselves, stop the sequence.

In order for a singularity to form, velocity has to blow up. In order for
that to happen, its acceleration would also have to blow up, which means its
jerk has to blow up, and so on, and so on:
\[
u,\qquad \partial_tu,\qquad \partial_t^2u,\qquad \ldots
\]
What you'd expect to see in this pattern, if a singularity were to form, as
you go up the time derivatives, you'd expect to see less spatial regularity,
not more.

The equation moves in the other direction. Its first three time derivatives
are
\[
\begin{aligned}
\partial_tu
={}&\nu\Delta u
 -(u\!\cdot\nabla)u
 -\nabla p,\\[3pt]
\partial_t^2u
={}&\nu\Delta\partial_tu
 -(\partial_tu\!\cdot\nabla)u
 -(u\!\cdot\nabla)\partial_tu
 -\nabla\partial_tp,\\[3pt]
\partial_t^3u
={}&\nu\Delta\partial_t^2u
 -(\partial_t^2u\!\cdot\nabla)u
 -2(\partial_tu\!\cdot\nabla)\partial_tu
 -(u\!\cdot\nabla)\partial_t^2u
 -\nabla\partial_t^2p.
\end{aligned}
\]
\begin{samepage}
The pattern is
\[
\partial_t^{\,n+1}u
=\nu\Delta\partial_t^{\,n}u
-\sum_{a=0}^{n}\binom na
  \bigl(\partial_t^{\,a}u\!\cdot\nabla\bigr)
  \partial_t^{\,n-a}u
-\nabla\partial_t^{\,n}p.
\]
\end{samepage}
\begin{samepage}
Set
\[
V_n:=\partial_t^{\,n}u,
\qquad
Q_n:=\partial_t^{\,n}p,
\]
Taking the divergence and using incompressibility gives the pressure at that
order:
\[
-\Delta Q_n
=\sum_{a=0}^{n}\binom na
  \partial_i\partial_j
  \bigl((V_a)_i(V_{n-a})_j\bigr).
\]
\end{samepage}
The line to keep in view is \(V_{n+1}=\nu\Delta V_n-\cdots\). Moving one step
up in time brings in two spatial derivatives of the response below it.
Transport and pressure stay attached, but viscosity supplies the link from
temporal order to spatial Sobolev order.

At the critical height, set
\[
\Lambda=(-\Delta)^{1/2},
\qquad
E_s(t):=\frac12\|\Lambda^su(t)\|_2^2,
\qquad
H(t):=E_{1/2}(t).
\]
If \(\mathcal P_s\) denotes the complete transport--pressure contribution at
height \(s\), the energy identity is
\[
E_s'=\mathcal P_s-2\nu E_{s+1}.
\]
Starting from \(H=E_{1/2}\) and repeating that identity gives
\[
H^{(n)}
=(-2\nu)^nE_{n+1/2}
 +\sum_{j=0}^{n-1}(-2\nu)^j
   \partial_t^{\,n-1-j}\mathcal P_{j+1/2}.
\]
The first term has reached spatial height \(n+\tfrac12\); the sum carries every
transport--pressure term generated on the way. This is the opposite of the
blow-up picture. Rising through that temporal sequence was supposed to reveal
less spatial regularity. The equation exposes higher Sobolev height instead. As
you do that, you get higher order Sobolev embeddings:
\[
s>k+\frac32
\quad\Longrightarrow\quad
H^s(\mathbb R^3)\hookrightarrow C^k(\mathbb R^3),
\qquad
\bigcap_{s<\infty}H^s(\mathbb R^3)
=H^\infty(\mathbb R^3)
\hookrightarrow C^\infty(\mathbb R^3).
\]

Seeing a higher energy in the identity is not enough. It has to remain finite
all the way to \(T_*\). Fix smooth divergence-free finite-energy data \(u_0\),
viscosity \(\nu>0\), and its maximal pressure-complete solution \((u,p)\) on
\([0,T_*)\). Assume \(T_*<\infty\). For finite \(N\) and \(M\), collect the
adjacent rows in
\[
\mathfrak R_{N,M}(u;T)
:=\sum_{n=0}^{N}\left(
  \|V_n\|_{L^2(0,T;H^{M+1})}^2
 +\|V_{n+1}\|_{L^2(0,T;H^{M-1})}^2
\right).
\]
Adding the corresponding \(Q_n\) and \(\nabla Q_n\) coordinates gives the
finite pressure-complete rectangle \(\mathcal R_{N,M}[u_0;T]\). The
manuscript's whole-terminal theorem bounds each adjacent pair on the full
interval:
\[
V_n\in L^2(0,T_*;H^{M+1}),
\qquad
V_{n+1}\in L^2(0,T_*;H^{M-1}),
\qquad 0\le n\le N.
\]
Restriction to \((0,T)\) can only decrease those norms, and there are only
finitely many rows, so
\[
\sup_{0<T<T_*}\mathfrak R_{N,M}(u;T)<\infty.
\]
\begin{samepage}
The bound may depend on \(N\) and \(M\). No uniform constant over all orders is
needed. At \(n=0\), for every finite \(m\), the adjacent pair gives
\[
u\in L^2(0,T_*;H^{m+1}),
\qquad
\partial_tu\in L^2(0,T_*;H^{m-1}),
\]
The Lions--Magenes trace theorem applied to
\(H^{m+1}\hookrightarrow H^m\hookrightarrow H^{m-1}\), gives a unique
trace at the endpoint:
\[
u\in C([0,T_*];H^m)
\qquad\text{for every finite }m.
\]
\end{samepage}

\begin{samepage}
Increasing \(N\) or \(M\) does not create a new solution. Every rectangle is a
finite part of the maximal history from \(u_0\), so overlapping rectangles
agree on every shared coordinate. Their projective limit is
\[
\mathcal I[u_0]
:=\varprojlim_{N,M<\infty}\mathcal R_{N,M}[u_0;T_*],
\qquad
u\in\bigcap_{r,m<\infty}H^r(0,T_*;H^m(\mathbb R^3)).
\]
\end{samepage}
The intersection is not one bound over infinitely many orders. Each coordinate
was bounded at a finite \((N,M)\); compatibility says that all of those finite
bounds belong to one velocity--pressure history.

The zeroth temporal row, \(\mathscr S_0:=\pi_0\mathcal I[u_0]\), now has one
terminal value at every finite spatial height:
\[
u_*:=u(T_*)
\in\bigcap_{m<\infty}H^m_\sigma(\mathbb R^3)
=H^\infty_\sigma(\mathbb R^3)
\hookrightarrow C^\infty(\mathbb R^3).
\]
The recurrence and pressure equation recover every finite coordinate of the
terminal jet from \(u_*\). The uniform \(H^3\) row gives a local existence time
\(\delta_0>0\) that works for \(u_*\) and for every sufficiently late slice
\(u(t_0)\). Choose \(t_0<T_*\) with \(T_*-t_0<\delta_0/2\), and restart from
\(u(t_0)\). Uniqueness makes this restart identical to the original solution
on \([t_0,T_*)\), so it reaches \(T_*\) with value \(u_*\). To the right of
\(T_*\), uniqueness identifies it with the solution restarted from \(u_*\).
The two restarts join into one smooth pressure-complete solution beyond
\(T_*\), contradicting maximality. Thus \(T_*=\infty\).

Once \(V_r\in C_tH_x^m\) is available for every finite \(r\) and \(m\), any
finite downstream estimate comes from choosing
\[
m>k+\frac32-\frac3p,
\qquad
2\le p\le\infty,
\]
which gives, on every finite time interval,
\[
V_r\in C_tH_x^m
\hookrightarrow L_t^qW_x^{k,p},
\qquad 1\le q\le\infty.
\]

\[
u_0
\longrightarrow \{(V_n,Q_n)\}_{n\ge0}
\longrightarrow \{\mathcal R_{N,M}[u_0]\}_{N,M<\infty}
\longrightarrow \mathcal I[u_0]
\longrightarrow \mathscr S_0
\longrightarrow u_*\in H^\infty
\longrightarrow T_*=\infty.
\]
